\chapter{\IfLanguageName{dutch}{Stand van zaken}{State of the art}}%
\label{ch:stand-van-zaken}

% Tip: Begin elk hoofdstuk met een paragraaf inleiding die beschrijft hoe
% dit hoofdstuk past binnen het geheel van de bachelorproef. Geef in het
% bijzonder aan wat de link is met het vorige en volgende hoofdstuk.

% Pas na deze inleidende paragraaf komt de eerste sectiehoofding.

Dit hoofdstuk bouwt verder op de inleiding en verdiept zich in de bestaande literatuur rond verkoopvoorspelling en datagedreven personeelsplanning. Waar in het vorige hoofdstuk het probleem en de onderzoeksvragen werden gedefinieerd, wordt hier het wetenschappelijke kader geschetst waarbinnen deze bachelorproef zich situeert. Het doel is om inzicht te geven in de huidige stand van zaken binnen het domein van retail forecasting en machine learning, zodat de gemaakte methodologische keuzes in het volgende hoofdstuk onderbouwd zijn.

\subsection*{1. Vergelijken modelprestaties}
De literatuur toont een duidelijk onderscheid in prestaties tussen traditionele lineaire modellen en moderne ensemble-methoden. \cite{GangulyMukherjee2024} vergeleek verschillende machine-learningmodellen voor retailverkoopvoorspelling en stelde concrete verschillen vast in prestatie: Random Forest behaalde een R² van 0.945 met een RMSE van 282.07, terwijl lineaire regressie slechts een R² van 0.531 bereikte. Deze resultaten tonen aan dat ensemble-methoden aanzienlijk sterker presteren bij het modelleren van verkooppatronen.

Gradient Boosting behoort ook tot de best presterende modellen. \cite{GangulyMukherjee2024} rapporteert een R² van 0.942, terwijl \cite{Kang2023} aantoont dat Gradient Boosting in hun studie marginaal beter presteerde dan Random Forest. Daarnaast toont \cite{Easter2024} aan dat XGBoost robuuster presteert dan klassieke tijdreeksmodellen zoals SARIMA bij dagelijkse bakkerijvraagvoorspellingen, met MAE-waarden tussen 0.0041 en 2.0978 en RMSE-waarden tussen 0.1112 en 3.5450. Deze consistente prestaties op verschillende datasets wijzen erop dat ensemble-methoden geschikt zijn voor het voorspellen van retailverkopen.

\subsection*{2. Rol van operationele en contextuele data}
\cite{Geboers2012} onderzocht in zijn case study bij Bakkersland B.V. de impact van contextuele variabelen op verkoopvoorspellingen voor de broodindustrie. Het concludeert dat het toevoegen van contextuele data, zoals weersinvloeden, feestdagen en lokale evenementen, de nauwkeurigheid significant verbetert ten opzichte van modellen die enkel historische verkoopdata gebruiken.

\cite{Easter2024} benadrukt dat de preprocessing van data cruciaal is voor modelprestaties. Het smoothing van tijdsreeksen vermindert de impact van outliers en volatiliteit, wat zowel traditionele als moderne modellen ten goede komt. Voor bakkerijen is deze combinatie van zorgvuldige data-voorbereiding en contextuele variabelen essentieel om robuuste voorspellingen te bekomen.

\subsection*{3. Van voorspelling naar personeelsplanning}
\cite{Geboers2012} beschrijft in zijn case study hoe verkoopvoorspellingen kunnen worden vertaald naar productieplanning. Voor personeelsplanning vereist dit het definiëren van capaciteitsregels: hoeveel klanten kan één medewerker per tijdseenheid bedienen, wat is de minimale bezetting ongeacht voorspelde drukte, en welke buffer is nodig voor onzekerheid?

\cite{Subekti2025} demonstreert dat convolutionele neurale netwerken uurbasis voorspellingen kunnen maken voor individuele bakkerijproducten. Detailniveau maakt een nauwkeurigere personeelsplanning mogelijk door verkoopvolumes per tijdsblok direct om te zetten naar arbeidsbehoefte. Dergelijke benaderingen creëren een coherente keten van data naar operationele beslissingen.

\subsection*{4. Conclusie}
De literatuur toont consistent aan dat ensemble-methoden zoals Random Forest, Gradient Boosting en XGBoost significant beter presteren dan traditionele lineaire modellen voor het voorspellen van verkoop. Deep learning-technieken zoals convolutionele neurale netwerken bieden aanvullende mogelijkheden voor fijnmazige, uurbasis voorspellingen.

Voor bakkerijen is het integreren van zowel operationele data als contextuele data cruciaal voor een nauwkeurige voorspelling. Een zorgvuldige data-preprocessing versterkt deze effecten verder. De combinatie van deze data en machine-learningmodellen stelt kleine ondernemingen in staat om de personeelsplanning te optimaliseren.

De uitdaging ligt in het praktisch implementeren van deze technologieën. Een proof-of-concept moet niet alleen robuust zijn, maar ook praktisch toepasbaar en gebruiksvriendelijk voor de eindgebruiker.
